\scanpage{14}
vu que $\deleted{\mathbb F_3}$ l'on \unsure{a}
\[
\mathbb F_3^*=\{\pm1\}\qquad\big).
\]

\deleted{Ainsi} \inserted{Notons que} la donnée de $\deleted{\{}(\omega_i)_{i\in\tau}\deleted{\}}$ permet ainsi de réconstituer \inserted{(à isom. can. près)} la famille $(D_i)_{i\in\tau}$ de droites \inserted{vectorielles sur $\mathbb F_3$} (via $D_i=\omega_i\amalg\{0\}\ \ldots$), donc le vectoriel $V=\prod_{i\in\tau}D_i$ sur $\mathbb F_3$, donc aussi \inserted{le plan vectoriel} $V=\prod_{i\in\tau}D_i$ avec la paire de droites indexée par $\tau$, d'où la paire complémentaire $(D_i)_{i\in\tau'}$. La \deleted{\unsure{On a|trouve|Donc} ainsi en fait} donnée d'une famille $(\omega_i)_{\unsure{v|i}\in\tau}$ d'ens à deux éléments, indexée par $\tau$, équivaut à celle d'un \uline{carré} \inserted{comb.} $Q$, admettant $\tau$ comme ens des diagonales. On trouve \inserted{\unsure{ainsi}} \deleted{\unsure{donc|d'une|une}} équivalence

\diagram{\begin{tikzcd}[ampersand replacement=\&,column sep=large]
{\llap{\text{(9)}\quad}\begin{array}[t]{@{}l@{}}
\text{Carrés combinatoires}\\Q
\end{array}} \arrow[r,leftrightarrow,"\approx"] \&
{\begin{array}[t]{@{}l@{}}
\text{plans vectoriels }\inserted{V}\text{ sur }\mathbb F_3,\\
\text{munis d'une paire}\\
\text{de droites vectorielles}\\
(D_{\unsure{v|i}})_{i\in\tau}
\end{array}}
\end{tikzcd}}

dans cette équivalence, donc, \deleted{$\tau$ correspond à} \inserted{on a}

\deleted{\unsure{diag|Diag $Q$}}
\[
\left\{\begin{aligned}
\tau&\simeq\mathrm{diag}_{Q} &&\text{et}\quad D_i^*\simeq d_i\quad\text{pour }i\in\tau\\
S_Q&=\bigcup_{i\in\tau}D_i^*=\coprod_{i\in\tau}D_i^*\\
\inserted{A_Q}&\inserted{=\prod_{i\in\tau}D_i^*\simeq V^*-\bigcup_{\unsure{i\in\tau|i\in T}}D_i^*}\;
 \tikz[baseline=-.25ex]{\draw (0,0)--(.10,.16)--(.22,.42) (.10,.16)--(.21,0);}
\end{aligned}\right.
\tag*{(10)}
\]
\margin{$\tikz[baseline=-.25ex]{\draw (0,0)--(.10,.16)--(.22,.42) (.10,.16)--(.21,0);}
=\bigcup_{i\in\tau'} D_i^*$}
où $d_i$ est l'ensemble des deux sommets de $Q$ qui constituent la diagonale $i$
% archive sheet number: 13
